\subsubsection{Word2Vec Facts}\label{sec:word2vec-facts}

The introduction of \textbf{Word2Vec} marked a turning point in Natural Language Processing because it demonstrated that \textbf{simpler neural architectures,\break when trained on massive datasets, can outperform much more complex language models}. Rather than focusing on building deeper neural networks, Mikolov et al. concentrated on making training extremely efficient, allowing the model to learn from billions of words.

\highspace
\begin{flushleft}
    \textcolor{Green3}{\faIcon{tachometer-alt} \textbf{Training Speed}}
\end{flushleft}
One of the greatest strengths of Word2Vec is its \textbf{training efficiency}. Compared to Bengio's Neural Network Language Model (NNLM), Word2Vec \hl{removes the hidden layer} and greatly simplifies the computations performed for each training example. As a result, the computational complexity decreases from:
\begin{equation*}
    \mathcal{O}\left(nm + nmh + h\left|V\right|\right)
\end{equation*}
To approximately:
\begin{equation*}
    \mathcal{O}\left(nm + m\log\left|V\right|\right)
\end{equation*}
Where $n$ is the context size, $m$ is the embedding dimension, $V$ is the vocabulary size. The expensive hidden-layer computations disappear, \hl{leaving only the embedding lookup} and the \hl{output prediction}. This seemingly simple modification reduces the computational cost dramatically.

\highspace
The improvement is not merely theoretical. According to the original experiments reported by Mikolov, Word2Vec processes approximately $100,000$ to $5,000,000$ words \textbf{every second}. This made it feasible, for the first time, to train high-quality embeddings on web-scale corpora containing billions of words.

\highspace
\textcolor{Green3}{\faIcon{check-circle} \textbf{Scalability.}} Because of its computational efficiency, Word2Vec scales extremely well. Instead of being limited to relatively small corpora, it can exploit: millions of documents, billions of words, vocabularies containing millions of unique terms. As the amount of available text increases, the learned embeddings generally improve because the model observes many more contexts for each word.

\highspace
Therefore, \hl{rather than increasing the complexity of the neural network,\break Word2Vec increases the amount of data used for learning.} Modern NLP follows exactly this philosophy: \textbf{large datasets often compensate for simpler architectures}.

\highspace
\textcolor{Green3}{\faIcon[regular]{bookmark} \textbf{Large Corpora.}} The scalability of Word2Vec is evident from the experiments reported in the original paper. The model was trained on the \textbf{Google News} corpus containing approximately \textbf{6 billion words} and a vocabulary of roughly $\left|V\right| \approx 10^{6}$ unique words. Despite this enormous dataset, training required only:
\begin{itemize}
    \item \textbf{CBOW}: about \textbf{2 days} on \textbf{140 CPU cores}.
    \item \textbf{Skip-gram}: about \textbf{2.5 days} on \textbf{125 CPU cores}.
\end{itemize}
For comparison, Bengio's original Neural Network Language Model required: \textbf{14 days}, \textbf{180 CPU cores}, and an embedding size of only \textbf{100 dimensions}. Furthermore, the Word2Vec was trained using a much larger embedding size of 1000 dimensions, which would have been computationally prohibitive for the original NNLM.

\highspace
\textcolor{Green3}{\faIcon{check-circle} \textbf{Accuracy Improvements.}} Perhaps the most surprising result was that Word2Vec was \textbf{not only faster}, but also \textbf{more accurate}.
\begin{center}
    \begin{tabular}{@{} l l @{}}
        \toprule
        \textbf{Model} & \textbf{Accuracy} \\
        \midrule
        Best Neural Network Language Model (Bengio) & 12.3\% \\
        Word2Vec (Skip-Gram) & 53.3\% \\
        \bottomrule
    \end{tabular}
\end{center}
\noindent
Thus, Word2Vec achieved both \textbf{dramatically faster training} and \textbf{substantially higher predictive performance}, demonstrating that learning high-quality word embeddings from massive corpora can be more beneficial than using a deeper, more computationally expensive language model.

\highspace
In general, Word2Vec's impact extends far beyond its original architecture. Its key contributions include:
\begin{itemize}
    \item Demonstrating that \textbf{distributional semantics} (``\emph{a word is known by the company it keeps}'') can be learned automatically from raw text;
    \item Showing that \textbf{simple shallow networks} can outperform deeper models when trained on enough data;
    \item Producing embeddings that capture rich semantic and syntactic relationships;
    \item Making large-scale unsupervised representation learning practical.
\end{itemize}
Perhaps most importantly, Word2Vec shifted the community's attention from \textbf{predicting the next word} to \textbf{learning high-quality vector representations}. These embeddings became the foundation for numerous subsequent NLP models, including GloVe, FastText, ELMo, BERT, and many modern transformer-based language models.